%======================================================================
% CONFIGURACIÓN
%======================================================================

\chapter{Configuración}

Este capítulo describe las opciones de configuración disponibles en Sandra UI Consola para personalizar el comportamiento del sistema según las necesidades de cada organización.

\section{Configuración General}

La configuración general permite ajustar los parámetros básicos del sistema:

\begin{itemize}
    \item Nombre de la instancia
    \item Zona horaria
    \item Idioma de la interfaz
    \item Formato de fecha y hora
    \item Configuración de notificaciones
\end{itemize}

\section{Configuración de Conexión}

La consola se comunica con Sandra Server mediante múltiples protocolos. Los parámetros de conexión incluyen:

\begin{itemize}
    \item \textbf{Host}: Dirección del servidor Sandra
    \item \textbf{Puerto}: Número de puerto para la conexión
    \item \textbf{Protocolo}: gRPC, REST o WebSocket
    \item \textbf{Timeout}: Tiempo máximo de espera para respuestas
    \item \textbf{Retry}: Configuración de reintentos en caso de fallo
\end{itemize}

La comunicación multicanal permite optimizar cada tipo de operación según sus requisitos de latencia y fiabilidad:

\begin{itemize}
\item \textbf{gRPC}: Streaming de datos en tiempo real y operaciones de alta velocidad.
\item \textbf{REST}: Operaciones CRUD estándar y consultas.
\item \textbf{WebSockets}: Notificaciones push y actualizaciones en vivo.
\end{itemize}

\section{Parámetros Avanzados}

Los parámetros avanzados permiten un control más granular del sistema:

\begin{itemize}
    \item Configuración de memoria caché
    \item Límites de conexiones simultáneas
    \item Configuración de logs
    \item Parámetros de seguridad
    \item Configuración de alta disponibilidad
\end{itemize}

\section{Certificados SSL}

La gestión de certificados SSL es fundamental para la seguridad de las comunicaciones:

\begin{itemize}
    \item Importar certificados firmados por CA
    \item Generar certificados autofirmados
    \item Validar certificado antes de conectar
    \item Renovación automática de certificados
\end{itemize}

\clearpage
